\documentclass[a4paper]{article}

\usepackage[spanish]{babel}
\usepackage{graphicx}
\usepackage{amsmath, amssymb}
\usepackage[margin=2cm]{geometry}
\usepackage{fancyhdr}
\usepackage{enumerate}
\usepackage[shortlabels]{enumitem}
\usepackage{parskip}
\usepackage[most]{tcolorbox}
\usepackage[hidelinks]{hyperref}
\usepackage{float}
\usepackage{booktabs}



% cabecera
\pagestyle{fancy}
\fancyhead[l]{Fisica matemática I}
\fancyhead[c]{Problemas sección 1.6}
\fancyhead[r]{\today}
\fancyfoot[c]{\thepage}
\renewcommand{\headrulewidth}{0.2pt} % linea horizontal


\begin{document}

\section*{Solución: Cinemática en Coordenadas Cilíndricas (Página 18)}

\subsection*{Ítem 1: Derivadas temporales de los vectores unitarios}
\textbf{Enunciado:} Demuestre que $\dot{\hat{e}}_\rho = \hat{e}_\phi \dot{\phi}, \quad \dot{\hat{e}}_\phi = -\hat{e}_\rho \dot{\phi}, \quad \dot{\hat{e}}_z = 0$.

\begin{itemize}
    \item \textbf{Expresión en cartesianas:} Partimos de la definición de los vectores unitarios cilíndricos en términos de la base fija $\{\hat{i}, \hat{j}, \hat{k}\}$ (ver fragmento):
    \begin{align*}
        \hat{e}_\rho &= \hat{i} \cos\phi + \hat{j} \sin\phi \\
        \hat{e}_\phi &= -\hat{i} \sin\phi + \hat{j} \cos\phi \\
        \hat{e}_z &= \hat{k}
    \end{align*}
    
    \item \textbf{Derivada de $\hat{e}_\rho$:} Aplicando la regla de la cadena ($\frac{d}{dt} = \dot{\phi} \frac{d}{d\phi}$):
    \begin{equation*}
        \dot{\hat{e}}_\rho = \frac{d}{dt}(\hat{i} \cos\phi + \hat{j} \sin\phi) = (-\hat{i} \sin\phi + \hat{j} \cos\phi) \dot{\phi} = \hat{e}_\phi \dot{\phi} \text{}
    \end{equation*}
    
    \item \textbf{Derivada de $\hat{e}_\phi$:}
    \begin{equation*}
        \dot{\hat{e}}_\phi = \frac{d}{dt}(-\hat{i} \sin\phi + \hat{j} \cos\phi) = (-\hat{i} \cos\phi - \hat{j} \sin\phi) \dot{\phi} = -(\hat{i} \cos\phi + \hat{j} \sin\phi) \dot{\phi} = -\hat{e}_\rho \dot{\phi} \text{}
    \end{equation*}
    
    \item \textbf{Derivada de $\hat{e}_z$:} Como $\hat{e}_z = \hat{k}$ es un vector constante en dirección y magnitud, su derivada es nula: $\dot{\hat{e}}_z = 0$.
\end{itemize}

\subsection*{Ítem 2: Vector Velocidad ($\mathbf{v}$)}
\textbf{Enunciado:} Demuestre que $\mathbf{v} = \dot{\rho}\hat{e}_\rho + \rho\dot{\phi}\hat{e}_\phi + \dot{z}\hat{e}_z$.

\begin{itemize}
    \item \textbf{Vector posición:} En coordenadas cilíndricas, el vector posición se escribe como (ver fragmento):
    \begin{equation*}
        \mathbf{r} = \rho \hat{e}_\rho + z \hat{e}_z
    \end{equation*}
    
    \item \textbf{Derivación temporal:} La velocidad es la derivada de $\mathbf{r}$ respecto al tiempo:
    \begin{equation*}
        \mathbf{v} = \frac{d}{dt}(\rho \hat{e}_\rho + z \hat{e}_z) = \dot{\rho} \hat{e}_\rho + \rho \dot{\hat{e}}_\rho + \dot{z} \hat{e}_z + z \dot{\hat{e}}_z
    \end{equation*}
    
    \item \textbf{Sustitución:} Usando los resultados del Ítem 1:
    \begin{equation*}
        \mathbf{v} = \dot{\rho} \hat{e}_\rho + \rho (\dot{\phi} \hat{e}_\phi) + \dot{z} \hat{e}_z + z(0) = \dot{\rho} \hat{e}_\rho + \rho \dot{\phi} \hat{e}_\phi + \dot{z} \hat{e}_z \text{}
    \end{equation*}
\end{itemize}

\subsection*{Ítem 3: Vector Aceleración ($\mathbf{a}$)}
\textbf{Enunciado:} Demuestre que $\mathbf{a} = (\ddot{\rho} - \rho\dot{\phi}^2)\hat{e}_\rho + (\rho\ddot{\phi} + 2\dot{\rho}\dot{\phi})\hat{e}_\phi + \ddot{z}\hat{e}_z$.

\begin{itemize}
    \item \textbf{Derivación de la velocidad:} La aceleración es la derivada de $\mathbf{v}$:
    \begin{equation*}
        \mathbf{a} = \frac{d}{dt} (\dot{\rho} \hat{e}_\rho + \rho \dot{\phi} \hat{e}_\phi + \dot{z} \hat{e}_z)
    \end{equation*}
    
    \item \textbf{Expansión por regla del producto:}
    \begin{equation*}
        \mathbf{a} = (\ddot{\rho} \hat{e}_\rho + \dot{\rho} \dot{\hat{e}}_\rho) + (\dot{\rho} \dot{\phi} \hat{e}_\phi + \rho \ddot{\phi} \hat{e}_\phi + \rho \dot{\phi} \dot{\hat{e}}_\phi) + \ddot{z} \hat{e}_z
    \end{equation*}
    
    \item \textbf{Sustitución de derivadas de unitarios:}
    \begin{equation*}
        \mathbf{a} = \ddot{\rho} \hat{e}_\rho + \dot{\rho} (\dot{\phi} \hat{e}_\phi) + \dot{\rho} \dot{\phi} \hat{e}_\phi + \rho \ddot{\phi} \hat{e}_\phi + \rho \dot{\phi} (-\dot{\phi} \hat{e}_\rho) + \ddot{z} \hat{e}_z
    \end{equation*}
    
    \item \textbf{Agrupación de componentes:}
    \begin{itemize}
        \item \textbf{Componente $\hat{e}_\rho$:} $\ddot{\rho} - \rho \dot{\phi}^2$
        \item \textbf{Componente $\hat{e}_\phi$:} $\dot{\rho} \dot{\phi} + \rho \ddot{\phi} + \dot{\rho} \dot{\phi} = \rho \ddot{\phi} + 2 \dot{\rho} \dot{\phi}$
        \item \textbf{Componente $\hat{e}_z$:} $\ddot{z}$
    \end{itemize}
    
    \item \textbf{Resultado final:}
    \begin{equation*}
        \mathbf{a} = (\ddot{\rho} - \rho\dot{\phi}^2)\hat{e}_\rho + (\rho\ddot{\phi} + 2\dot{\rho}\dot{\phi})\hat{e}_\phi + \ddot{z}\hat{e}_z \text{}
    \end{equation*}
\end{itemize}

\section*{Solución: Cinemática en Coordenadas Esféricas (Página 18)}

Para resolver estos ítems, utilizaremos las derivadas parciales de los vectores unitarios esféricos respecto a sus coordenadas, calculadas previamente en la página 9 del texto \cite{50, 51}.

\subsection*{Ítem 1: Derivadas temporales de los vectores unitarios}
\textbf{Enunciado:} Demuestre que:
\begin{align*}
    \dot{\hat{e}}_r &= \hat{e}_\theta \dot{\theta} + \hat{e}_\phi \dot{\phi} \text{ sen } \theta \\
    \dot{\hat{e}}_\theta &= -\hat{e}_r \dot{\theta} + \hat{e}_\phi \dot{\phi} \text{ cos } \theta \\
    \dot{\hat{e}}_\phi &= -\hat{e}_\theta \dot{\phi} \text{ cos } \theta - \hat{e}_r \dot{\phi} \text{ sen } \theta
\end{align*}

\begin{itemize}
    \item \textbf{Procedimiento:} Aplicamos la regla de la cadena para la derivada temporal: $\dot{\hat{e}}_i = \frac{\partial \hat{e}_i}{\partial \theta}\dot{\theta} + \frac{\partial \hat{e}_i}{\partial \phi}\dot{\phi}$ (las derivadas respecto a $r$ son nulas \cite{50}).
    \item \textbf{Para $\hat{e}_r$:} Usando $\frac{\partial \hat{e}_r}{\partial \theta} = \hat{e}_\theta$ y $\frac{\partial \hat{e}_r}{\partial \phi} = \hat{e}_\phi \text{ sen } \theta$ \cite{50}:
    \begin{equation*}
        \dot{\hat{e}}_r = (\hat{e}_\theta) \dot{\theta} + (\hat{e}_\phi \text{ sen } \theta) \dot{\phi} = \hat{e}_\theta \dot{\theta} + \hat{e}_\phi \dot{\phi} \text{ sen } \theta \text{ \cite{72}}
    \end{equation*}
    \item \textbf{Para $\hat{e}_\theta$:} Usando $\frac{\partial \hat{e}_\theta}{\partial \theta} = -\hat{e}_r$ y $\frac{\partial \hat{e}_\theta}{\partial \phi} = \hat{e}_\phi \text{ cos } \theta$ \cite{50, 51}:
    \begin{equation*}
        \dot{\hat{e}}_\theta = (-\hat{e}_r) \dot{\theta} + (\hat{e}_\phi \text{ cos } \theta) \dot{\phi} = -\hat{e}_r \dot{\theta} + \hat{e}_\phi \dot{\phi} \text{ cos } \theta \text{ \cite{72}}
    \end{equation*}
    \item \textbf{Para $\hat{e}_\phi$:} Usando $\frac{\partial \hat{e}_\phi}{\partial \phi} = -\hat{e}_r \text{ sen } \theta - \hat{e}_\theta \text{ cos } \theta$ \cite{51}:
    \begin{equation*}
        \dot{\hat{e}}_\phi = \frac{\partial \hat{e}_\phi}{\partial \phi} \dot{\phi} = (-\hat{e}_r \text{ sen } \theta - \hat{e}_\theta \text{ cos } \theta) \dot{\phi} = -\hat{e}_\theta \dot{\phi} \text{ cos } \theta - \hat{e}_r \dot{\phi} \text{ sen } \theta \text{ \cite{72}}
    \end{equation*}
\end{itemize}

\subsection*{Ítem 2: Vector Velocidad ($\mathbf{v}$)}
\textbf{Enunciado:} Demuestre que $\mathbf{v} = \hat{e}_r \dot{r} + \hat{e}_\theta r\dot{\theta} + \hat{e}_\phi r\dot{\phi} \text{ sen } \theta$.

\begin{itemize}
    \item \textbf{Vector posición:} En esféricas, $\mathbf{r} = r \hat{e}_r$ \cite{47}.
    \item \textbf{Derivación:} 
    \begin{equation*}
        \mathbf{v} = \frac{d}{dt}(r \hat{e}_r) = \dot{r} \hat{e}_r + r \dot{\hat{e}}_r
    \end{equation*}
    \item \textbf{Sustitución:} Usando el resultado del Ítem 1 para $\dot{\hat{e}}_r$:
    \begin{equation*}
        \mathbf{v} = \dot{r} \hat{e}_r + r (\hat{e}_\theta \dot{\theta} + \hat{e}_\phi \dot{\phi} \text{ sen } \theta) = \hat{e}_r \dot{r} + \hat{e}_\theta r\dot{\theta} + \hat{e}_\phi r\dot{\phi} \text{ sen } \theta \text{ \cite{72}}
    \end{equation*}
\end{itemize}

\subsection*{Ítem 3: Vector Aceleración ($\mathbf{a}$)}
\textbf{Enunciado:} Demuestre que:
\begin{align*}
    \mathbf{a} = &\hat{e}_r(\ddot{r} - r\dot{\theta}^2 - r\dot{\phi}^2 \text{ sen}^2 \theta) \\
    &+ \hat{e}_\theta(2\dot{r}\dot{\theta} + r\ddot{\theta} - r\dot{\phi}^2 \text{ sen } \theta \text{ cos } \theta) \\
    &+ \hat{e}_\phi(2\dot{r}\dot{\phi} \text{ sen } \theta + r\ddot{\phi} \text{ sen } \theta + 2r\dot{\theta}\dot{\phi} \text{ cos } \theta)
\end{align*}

\begin{itemize}
    \item \textbf{Procedimiento:} Derivamos término a término el vector velocidad: $\mathbf{a} = \frac{d}{dt} (\mathbf{v}_r + \mathbf{v}_\theta + \mathbf{v}_\phi)$.
    \item \textbf{Derivada de cada componente:}
    \begin{enumerate}
        \item $\frac{d}{dt}(\dot{r}\hat{e}_r) = \ddot{r}\hat{e}_r + \dot{r}(\hat{e}_\theta \dot{\theta} + \hat{e}_\phi \dot{\phi} \text{ sen } \theta)$
        \item $\frac{d}{dt}(r\dot{\theta}\hat{e}_\theta) = \dot{r}\dot{\theta}\hat{e}_\theta + r\ddot{\theta}\hat{e}_\theta + r\dot{\theta}(-\hat{e}_r \dot{\theta} + \hat{e}_\phi \dot{\phi} \text{ cos } \theta)$
        \item $\frac{d}{dt}(r\dot{\phi} \text{ sen } \theta \hat{e}_\phi) = (\dot{r}\dot{\phi} \text{ sen } \theta + r\ddot{\phi} \text{ sen } \theta + r\dot{\phi}\dot{\theta} \text{ cos } \theta)\hat{e}_\phi + (r\dot{\phi} \text{ sen } \theta)\dot{\hat{e}}_\phi$
    \end{enumerate}
    \item \textbf{Agrupación por componentes:}
    \begin{itemize}
        \item \textbf{En $\hat{e}_r$:} $\ddot{r} - r\dot{\theta}^2 - (r\dot{\phi} \text{ sen } \theta)(\dot{\phi} \text{ sen } \theta) = \ddot{r} - r\dot{\theta}^2 - r\dot{\phi}^2 \text{ sen}^2 \theta$
        \item \textbf{En $\hat{e}_\theta$:} $\dot{r}\dot{\theta} + \dot{r}\dot{\theta} + r\ddot{\theta} - (r\dot{\phi} \text{ sen } \theta)(\dot{\phi} \text{ cos } \theta) = 2\dot{r}\dot{\theta} + r\ddot{\theta} - r\dot{\phi}^2 \text{ sen } \theta \text{ cos } \theta$
        \item \textbf{En $\hat{e}_\phi$:} $\dot{r}\dot{\phi} \text{ sen } \theta + \dot{r}\dot{\phi} \text{ sen } \theta + r\ddot{\phi} \text{ sen } \theta + r\dot{\theta}\dot{\phi} \text{ cos } \theta + r\dot{\phi}\dot{\theta} \text{ cos } \theta$ \\
        $= 2\dot{r}\dot{\phi} \text{ sen } \theta + r\ddot{\phi} \text{ sen } \theta + 2r\dot{\theta}\dot{\phi} \text{ cos } \theta$
    \end{itemize}
    \item \textbf{Conclusión:} Al unir las componentes se obtiene la expresión final de la aceleración \cite{72}.
\end{itemize}

\section*{Solución: Problemas de la Página 18 (Geometría e Identidades)}

\subsection*{Parte 1: Factores de escala y elementos diferenciales en coordenadas cilíndricas}
\textbf{Enunciado:} Calcular los factores de escala y los elementos de línea, superficie y volumen en coordenadas cilíndricas $(\rho, \phi, z)$.

\begin{itemize}
    \item \textbf{Paso 1: Relaciones de transformación.}
    Las coordenadas cartesianas se expresan en términos de las cilíndricas como \cite{34}:
    \begin{equation*}
        x = \rho \cos\phi, \quad y = \rho \sin\phi, \quad z = z
    \end{equation*}
    El vector posición es: $\mathbf{r} = (\rho \cos\phi)\hat{i} + (\rho \sin\phi)\hat{j} + z\hat{k}$.

    \item \textbf{Paso 2: Factores de escala ($h_i$).}
    Se calculan mediante la magnitud de los vectores tangentes $\frac{\partial \mathbf{r}}{\partial u_i}$ \cite{43}:
    \begin{align*}
        h_\rho &= \left| \frac{\partial \mathbf{r}}{\partial \rho} \right| = | \cos\phi \hat{i} + \sin\phi \hat{j} | = \sqrt{\cos^2\phi + \sin^2\phi} = 1 \\
        h_\phi &= \left| \frac{\partial \mathbf{r}}{\partial \phi} \right| = | -\rho \sin\phi \hat{i} + \rho \cos\phi \hat{j} | = \sqrt{\rho^2(\sin^2\phi + \cos^2\phi)} = \rho \\
        h_z &= \left| \frac{\partial \mathbf{r}}{\partial z} \right| = | \hat{k} | = 1
    \end{align*}
    Por lo tanto: $h_\rho = 1, h_\phi = \rho, h_z = 1$ \cite{73}.

    \item \textbf{Paso 3: Elemento diferencial de línea ($d\mathbf{r}$).}
    Utilizando $d\mathbf{l}_i = h_i \hat{e}_i du_i$ \cite{69}:
    \begin{equation*}
        d\mathbf{r} = \hat{e}_\rho d\rho + \hat{e}_\phi \rho d\phi + \hat{e}_z dz
    \end{equation*}

    \item \textbf{Paso 4: Elementos diferenciales de superficie ($d\mathbf{S}_i$).}
    Definidos como $d\mathbf{S}_i = \hat{e}_i (h_j h_k du_j du_k)$ \cite{70}:
    \begin{align*}
        d\mathbf{S}_\rho &= \hat{e}_\rho (h_\phi h_z d\phi dz) = \hat{e}_\rho (\rho d\phi dz) \\
        d\mathbf{S}_\phi &= \hat{e}_\phi (h_\rho h_z d\rho dz) = \hat{e}_\phi (d\rho dz) \\
        d\mathbf{S}_z &= \hat{e}_z (h_\rho h_\phi d\rho d\phi) = \hat{e}_z (\rho d\rho d\phi)
    \end{align*}

    \item \textbf{Paso 5: Elemento diferencial de volumen ($dV$).}
    Se define como $dV = h_\rho h_\phi h_z d\rho d\phi dz$ \cite{71}:
    \begin{equation*}
        dV = (1)(\rho)(1) d\rho d\phi dz = \rho d\rho d\phi dz
    \end{equation*}
\end{itemize}

\subsection*{Parte 2: Demostración de Identidades Vectoriales}
\textbf{Enunciado:} Demostrar las identidades dadas utilizando el símbolo de Levi-Civita $\epsilon_{ijk}$.

\begin{enumerate}
    \item \textbf{Identidad:} $\mathbf{A} \times \mathbf{B} = \sum_{ijk} \hat{e}_i \epsilon_{ijk} A_j B_k$
    \begin{itemize}
        \item \textbf{Demostración:} Partimos de la expansión de los vectores en la base ortonormal: $\mathbf{A} = \sum_j A_j \hat{e}_j$ y $\mathbf{B} = \sum_k B_k \hat{e}_k$.
        \begin{equation*}
            \mathbf{A} \times \mathbf{B} = \left( \sum_j A_j \hat{e}_j \right) \times \left( \sum_k B_k \hat{e}_k \right) = \sum_{j,k} A_j B_k (\hat{e}_j \times \hat{e}_k)
        \end{equation*}
        Usando la definición del producto vectorial entre unitarios (Ec. 1.23): $\hat{e}_j \times \hat{e}_k = \sum_i \epsilon_{jki} \hat{e}_i$ \cite{71}.
        Debido a la simetría cíclica $\epsilon_{jki} = \epsilon_{ijk}$, obtenemos:
        \begin{equation*}
            \mathbf{A} \times \mathbf{B} = \sum_{i,j,k} A_j B_k \epsilon_{ijk} \hat{e}_i \text{ (Q.E.D.)}
        \end{equation*}
    \end{itemize}

    \item \textbf{Identidad:} $\mathbf{A} \cdot \mathbf{B} \times \mathbf{C} = \sum_{ijk} \epsilon_{ijk} A_i B_j C_k$
    \begin{itemize}
        \item \textbf{Demostración:} Sustituimos el resultado anterior para $\mathbf{B} \times \mathbf{C} = \sum_{l,j,k} \epsilon_{ljk} B_j C_k \hat{e}_l$.
        \begin{equation*}
            \mathbf{A} \cdot (\mathbf{B} \times \mathbf{C}) = \left( \sum_i A_i \hat{e}_i \right) \cdot \left( \sum_{l,j,k} \epsilon_{ljk} B_j C_k \hat{e}_l \right)
        \end{equation*}
        Como $\hat{e}_i \cdot \hat{e}_l = \delta_{il}$ \cite{42}:
        \begin{equation*}
            \mathbf{A} \cdot (\mathbf{B} \times \mathbf{C}) = \sum_{i,j,k} A_i B_j C_k \epsilon_{ijk} \text{ (Q.E.D.)}
        \end{equation*}
    \end{itemize}

    \item \textbf{Identidad:} $\mathbf{A} \cdot \mathbf{B} \times \mathbf{C} = \mathbf{B} \cdot \mathbf{C} \times \mathbf{A} = \mathbf{C} \cdot \mathbf{A} \times \mathbf{B}$
    \begin{itemize}
        \item \textbf{Demostración:} El símbolo de Levi-Civita es invariante bajo permutaciones cíclicas de sus índices: $\epsilon_{ijk} = \epsilon_{jki} = \epsilon_{kij}$ \cite{71}.
        \begin{align*}
            \mathbf{A} \cdot (\mathbf{B} \times \mathbf{C}) &= \sum \epsilon_{ijk} A_i B_j C_k \\
            \mathbf{B} \cdot (\mathbf{C} \times \mathbf{A}) &= \sum \epsilon_{jki} B_j C_k A_i = \sum \epsilon_{ijk} A_i B_j C_k \\
            \mathbf{C} \cdot (\mathbf{A} \times \mathbf{B}) &= \sum \epsilon_{kij} C_k A_i B_j = \sum \epsilon_{ijk} A_i B_j C_k
        \end{align*}
        Todos los términos son idénticos \cite{73}.
    \end{itemize}

    \item \textbf{Identidad:} $\mathbf{A} \cdot \mathbf{B} \times \mathbf{A} = 0$
    \begin{itemize}
        \item \textbf{Demostración:} Utilizando la forma de sumatoria:
        \begin{equation*}
            \mathbf{A} \cdot (\mathbf{B} \times \mathbf{A}) = \sum_{i,j,k} \epsilon_{ijk} A_i B_j A_k
        \end{equation*}
        El símbolo $\epsilon_{ijk}$ es totalmente antisimétrico. Si se repite un índice (en este caso $i$ y $k$ están asociados a las componentes del mismo vector $\mathbf{A}$), el término se anula: $\epsilon_{ijk} = -\epsilon_{kji}$. Para cada término con índices $(i,j,k)$ existe un término opuesto $(k,j,i)$ que lo cancela dado que $A_i A_k = A_k A_i$. Por definición de Levi-Civita, si dos índices son iguales, el valor es 0 \cite{71, 73}.
    \end{itemize}
\end{enumerate}


\section*{Demostraciones de la Página 19: Símbolo de Levi-Civita}

Para estas demostraciones, se parte de la identidad fundamental de contracción (ecuación 1.24 de las fuentes):
\begin{equation}
    \sum_{k=1}^{3} \epsilon_{ijk} \epsilon_{lmk} = \delta_{il}\delta_{jm} - \delta_{im}\delta_{jl} \text{,}
\end{equation}

\subsection*{1. Demostración de $\sum_{j,k=1}^3 \epsilon_{ijk} \epsilon_{ljk} = 2\delta_{il}$}
\begin{itemize}
    \item \textbf{Paso 1:} Tomamos la identidad de contracción y sumamos sobre el índice $j$ haciendo $m = j$,.
    \item \textbf{Paso 2:} La expresión se transforma en:
    \begin{equation}
        \sum_{j=1}^{3} \left( \sum_{k=1}^{3} \epsilon_{ijk} \epsilon_{ljk} \right) = \sum_{j=1}^{3} (\delta_{il}\delta_{jj} - \delta_{ij}\delta_{jl}) \text{,}
    \end{equation}
    \item \textbf{Paso 3:} Evaluamos los términos resultantes. En tres dimensiones, la traza de la delta de Kronecker es $\sum_{j=1}^3 \delta_{jj} = 1 + 1 + 1 = 3$,. Por la propiedad de selección de la delta, $\sum_j \delta_{ij}\delta_{jl} = \delta_{il}$,.
    \item \textbf{Paso 4:} Restamos los valores obtenidos: $3\delta_{il} - \delta_{il} = 2\delta_{il}$,.
\end{itemize}

\subsection*{2. Demostración de $\sum_{i,j,k=1}^3 \epsilon_{ijk} \epsilon_{ijk} = 6$}
\begin{itemize}
    \item \textbf{Paso 1:} Partimos del resultado de la demostración anterior: $\sum_{j,k} \epsilon_{ijk} \epsilon_{ljk} = 2\delta_{il}$,.
    \item \textbf{Paso 2:} Realizamos la suma sobre el índice restante $i$ haciendo $l = i$,.
    \item \textbf{Paso 3:} La sumatoria se reduce a:
    \begin{equation}
        \sum_{i=1}^{3} 2\delta_{ii} \text{,}
    \end{equation}
    \item \textbf{Paso 4:} Como se estableció previamente, $\sum_{i=1}^3 \delta_{ii} = 3$,. Por lo tanto, el resultado final es $2 \times 3 = 6$,.
\end{itemize}

\subsection*{3. Demostración de $\sum_{j=1}^3 \delta_{ij} \epsilon_{ijk} = 0$}
\begin{itemize}
    \item \textbf{Paso 1:} La delta de Kronecker $\delta_{ij}$ actúa como un operador de selección dentro de la suma, obligando a que el índice $j$ sea igual a $i$ para que el término no sea nulo,.
    \item \textbf{Paso 2:} Al aplicar esta selección sobre el símbolo de Levi-Civita, obtenemos:
    \begin{equation}
        \sum_{j=1}^{3} \delta_{ij} \epsilon_{ijk} = \epsilon_{iik} \text{,}
    \end{equation}
    \item \textbf{Paso 3:} Por definición, el símbolo de Levi-Civita es totalmente antisimétrico y su valor es $0$ si existen dos o más índices repetidos (en este caso, el primer y segundo índice son $i$),,.
    \item \textbf{Conclusión:} La expresión es idénticamente nula para cualquier valor de $k$,.
\end{itemize}


\section*{Solución a los Problemas de la Página 19: Identidades Vectoriales}


Para estas demostraciones es indispensable utilizar la **identidad (1.24)** de las fuentes, la cual establece la contracción del producto de dos símbolos de Levi-Civita:
\begin{equation}
    \sum_{k=1}^{3} \epsilon_{ijk} \epsilon_{lmk} = \delta_{il}\delta_{jm} - \delta_{im}\delta_{jl} \text{}
\end{equation}


\subsection*{1. Demostración de $(\mathbf{A} \times \mathbf{B}) \cdot (\mathbf{C} \times \mathbf{D}) = (\mathbf{A} \cdot \mathbf{C})(\mathbf{B} \cdot \mathbf{D}) - (\mathbf{A} \cdot \mathbf{D})(\mathbf{B} \cdot \mathbf{C})$}

\begin{itemize}
    \item \textbf{Paso 1: Expresión en componentes.}
    Utilizamos la definición del producto vectorial mediante el símbolo de Levi-Civita $\epsilon_{ijk}$. La componente $i$ de cada producto es:
    \begin{equation}
        (\mathbf{A} \times \mathbf{B})_i = \sum_{j,k} \epsilon_{ijk} A_j B_k \quad \text{y} \quad (\mathbf{C} \times \mathbf{D})_i = \sum_{l,m} \epsilon_{ilm} C_l D_m
    \end{equation}
    
    \item \textbf{Paso 2: Producto escalar.}
    El producto escalar se define como la suma del producto de sus componentes homólogas:
    \begin{equation}
        (\mathbf{A} \times \mathbf{B}) \cdot (\mathbf{C} \times \mathbf{D}) = \sum_i \left( \sum_{j,k} \epsilon_{ijk} A_j B_k \right) \left( \sum_{l,m} \epsilon_{ilm} C_l D_m \right)
    \end{equation}
    Reordenando los sumatorios:
    \begin{equation}
        = \sum_{j,k,l,m} A_j B_k C_l D_m \left( \sum_i \epsilon_{ijk} \epsilon_{ilm} \right)
    \end{equation}

    \item \textbf{Paso 3: Aplicación de la identidad de contracción (1.24).}
    Aprovechamos la propiedad cíclica del símbolo ($\epsilon_{ijk} = \epsilon_{jki}$) para alinear el índice de suma:
    \begin{equation}
        \sum_i \epsilon_{jki} \epsilon_{lmi} = \delta_{jl}\delta_{km} - \delta_{jm}\delta_{kl}
    \end{equation}
    Sustituyendo en la expresión anterior:
    \begin{equation}
        = \sum_{j,k,l,m} A_j B_k C_l D_m (\delta_{jl}\delta_{km} - \delta_{jm}\delta_{kl})
    \end{equation}

    \item \textbf{Paso 4: Simplificación con la Delta de Kronecker.}
    Distribuimos los términos y aplicamos la propiedad selectiva de $\delta_{ij}$:
    \begin{itemize}
        \item Primer término ($\delta_{jl}\delta_{km}$): obliga a $l=j$ y $m=k \implies (\sum_j A_j C_j) (\sum_k B_k D_k) = (\mathbf{A} \cdot \mathbf{C})(\mathbf{B} \cdot \mathbf{D})$.
        \item Segundo término ($\delta_{jm}\delta_{kl}$): obliga a $m=j$ y $l=k \implies (\sum_j A_j D_j) (\sum_k B_k C_k) = (\mathbf{A} \cdot \mathbf{D})(\mathbf{B} \cdot \mathbf{C})$.
    \end{itemize}

    \item \textbf{Resultado:} Restando ambos términos se obtiene la identidad buscada.
\end{itemize}

\subsection*{2. Demostración de $(\mathbf{A} \times \mathbf{B}) \cdot (\mathbf{A} \times \mathbf{B}) = A^2 B^2 - (\mathbf{A} \cdot \mathbf{B})^2$}

\begin{itemize}
    \item \textbf{Procedimiento:}
    Esta identidad es un caso particular de la anterior. Si hacemos que los vectores $\mathbf{C} = \mathbf{A}$ y $\mathbf{D} = \mathbf{B}$, sustituimos directamente en el resultado obtenido en el punto 1:
    \begin{equation}
        (\mathbf{A} \times \mathbf{B}) \cdot (\mathbf{A} \times \mathbf{B}) = (\mathbf{A} \cdot \mathbf{A})(\mathbf{B} \cdot \mathbf{B}) - (\mathbf{A} \cdot \mathbf{B})(\mathbf{B} \cdot \mathbf{A})
    \end{equation}

    \item \textbf{Paso final:}
    Considerando que el producto escalar de un vector consigo mismo es su módulo al cuadrado ($\mathbf{A} \cdot \mathbf{A} = A^2$) y que el producto escalar es conmutativo ($\mathbf{A} \cdot \mathbf{B} = \mathbf{B} \cdot \mathbf{A}$):
    \begin{equation}
        (\mathbf{A} \times \mathbf{B}) \cdot (\mathbf{A} \times \mathbf{B}) = A^2 B^2 - (\mathbf{A} \cdot \mathbf{B})^2 \text{ (Q.E.D.)}
    \end{equation}
\end{itemize}



\bibliographystyle{plainurl}
\bibliography{bibliografia} 
\end{document}
